\documentclass[10pt,a4paper]{amsart}
\usepackage[margin=19mm]{geometry}
\usepackage{amsmath,amssymb,mathtools}
\usepackage[scr=rsfs,cal=euler]{mathalfa}
\usepackage{xfrac}
\usepackage[unicode,hidelinks,bookmarks=false]{hyperref}
\newcommand{\ie}{i.e.\ }
\newcommand{\cf}{cf.\ }
\newcommand{\resp}{respectively\ }
\newcommand{\Subsection}{\subsection}
\providecommand{\nobreakdash}{\nobreak\hbox{-}}
\setlength{\emergencystretch}{2em}
\multlinegap0pt
\usepackage{amssymb}
\usepackage{xcolor}
\newtheorem{prop}{Proposition}
\newtheorem{claim}{Claim}
\def\Ind{\setbox0=\hbox{$x$}\kern\wd0\hbox to 0pt{\hss$\mid$\hss} \lower.9\ht0\hbox to 0pt{\hss$\smile$\hss}\kern\wd0}
\def\ind{\mathop{\mathpalette\Ind{}}}
\newcommand{\m}{\mathbb }
\def\span{\operatorname{span}}
\title{A counterexample to the stable forking conjecture}
\author{James Freitag, Scott Mutchnik}
\date{Source: arXiv:2609.00436v1 (CC BY 4.0)}
\begin{document}
\maketitle

\noindent\emph{Source and licence.} A counterexample to the stable forking conjecture. Version arXiv:2609.00436v1 (CC BY 4.0), distributed by arXiv under the Creative Commons Attribution 4.0 International licence, http://creativecommons.org/licenses/by/4.0/, as recorded in the arXiv licence metadata for this exact version and shown on the abstract page https://arxiv.org/abs/2609.00436v1. Full licence notice: \texttt{licenses/arxiv-2609.00436-CC-BY-4.0.txt}.\par\smallskip

\noindent\textit{Editorial scope.} Counted unit: the article's prop prop (source0.tex lines 677--679). The article's complete original proof of it is delivered with it (source0.tex lines 681--722, one \texttt{proof} environment of 42 lines). No mathematics is written, repaired or re-derived: the statement and the proof are the article's own lines, in the article's own notation, and no retained line is substituted or re-flowed.  No further source-local context is needed: every cross-reference of every retained line resolves inside the two delivered spans.  Nothing was omitted from any retained span, so the package carries no omission record.  No retained line contains a citation, so the wrapper carries no bibliography and no bibliography entry.  No retained line contains a figure environment, an image-inclusion command or drawing code, so no image is delivered and the package declares no asset dependency.  Editorial formatting only, in addition: an AMS \texttt{amsart} preamble that restates the article's own declarations and macros used by the retained lines, namely the environments \texttt{prop}, \texttt{claim} on separate counters, together with the standard packages those lines need.  The retained environments print in this document as Proposition 1, Claim 1 in order of appearance, whereas the article prints the same items with the numbering of its own full document.  The complete native source of the article as distributed in the arXiv e-print for this exact version is bundled unchanged as \texttt{sources/209027/source0.tex}.
\begin{prop}
    If $M \prec \mathbb{M}$, $A_{1} \ind_{M} B$, $A_{2} \ind_{M} C$, $C \ind_{M} B$, and $A_{1} \equiv_{M} A_{2}$, then there is some $A \ind_{M} BC$ with $A \equiv_{MB} A_{1}$ and $A \equiv_{MC} A_{2}$.
\end{prop}

\begin{proof}
    Suppose that we have closed sets $D, E, F \leq \m M$ with $F \subseteq D \cap E$ and $D \ind _F E$. Take a $Q_{ \Gamma (F)}$-basis of $V(F)$, call it $B$. Then take an extension of this basis, $B \cup B_D$ to be a basis of $V(D)$ over the ring $Q _{\Gamma (D)}$. Similarly, let $B \cup B_E$ be a basis of $V( E)$ over $Q _{\Gamma (E)}$.

   % \begin{claim}
        Then we claim that $B \cup B_E \cup B_D$ is an independent set over $Q_{\Gamma (\m M)}$, and that the closure $\mathrm{cl}(DE) = (\Gamma (D) \cup \Gamma (E), \span _{Q_{\Gamma (D) \cup \Gamma (E)}} (B \cup B_E  \cup B_D))$. But this is just as in the proof of base monotonicity.
   % \end{claim}
%\begin{proof} (of claim) 
   % We know that both $B \cup B_D$ and $B \cup B_E$ are independent over $Q_{ \Gamma (\m M)}$ by the fact that $D,E \leq \m M$. So, suppose for a moment that there is a $Q_{ \Gamma (\m M)}$ relation on $B \cup B_E \cup B_D$. Moving things around, this is an equality of two elements - one in $\span_ {Q( \Gamma ( \m M ))} (B \cup B_E)$ and one in $\span_ {Q( \Gamma ( \m M ))} (B \cup B_D)$. 

%We know because of the fact that $D \ind _F E $ that the intersection of these two modules is $\span_ {Q( \Gamma ( \m M ))} (B)$. But since $B$ is independent over $Q( \Gamma ( \m M ))$, we must have that the relation is identically zero. 

%Thus, we know that $(\Gamma(E) \cup \Gamma (D) , \span_ {Q( \Gamma ( \m M ))} (B \cup B_D \cup B_E))) \leq \m M$. So, we've proved the claim. 
%\end{proof}

By our assumptions and definition of $\ind$, we have 
$$\mathrm{cl}(MA_1) \ind_M \mathrm{cl} (MB_0) , \, \mathrm{cl}(MA_2) \ind _M \mathrm{cl}(MC), \, \mathrm{cl}(MC) \ind _M \mathrm{cl} (MB).$$ (Here, note that $\mathrm{cl}(M) = M$, because $M \prec \mathbb{M}$ and, by definition of $\leq$, this implies $M \leq \mathbb{M}$.) By our elementary equivalence assumption, we have an isomorphism $f: \mathrm{cl}(MA_1) \rightarrow \mathrm{cl}(MA_2)$ which fixes $M$ pointwise and sends $A_1 $ to $A_2$. Pick a $Q _ { \Gamma (M)}$ basis $B_M$ of $V(M)$. Extend this basis to bases: 
\begin{enumerate}
    \item $B_M \cup B_{A_1} $ of $V(\mathrm{cl}(MA_1))$ over $Q_{\Gamma (\mathrm{cl}(MA_1))}$
    \item $B_M \cup f(B_{A_1}) $ of $V(\mathrm{cl}(MA_2))$ over $Q_{\Gamma (\mathrm{cl}(MA_2))}$
    \item $B_M \cup B_0 $ of $V(\mathrm{cl}(MB))$ over $Q_{\Gamma (\mathrm{cl}(MB))}$
    \item $B_M \cup C_0 $ of $V(\mathrm{cl}(MC))$ over $Q_{\Gamma (\mathrm{cl}(MC))}$
\end{enumerate}
By the hypothesis; and our claim following from the arguments from the proof of base monotonicity, we have the following: $B_M \cup B_{A_1} \cup B_0 $, $B_M \cup f(B_{A_1}) \cup C_0 $ and $B_M \cup C_0 \cup B_0 $ are each $Q_{\Gamma (\m M)}$-independent sets. 

Pick an embedding of $\Gamma (\mathrm{cl}( MA_1))$ into $\Gamma(\m M)$ which fixes $\Gamma (M)$ pointwise. Call it $\sigma$. And make sure that we have the property that $\sigma (\Gamma ( \mathrm{cl}(MA_1)))$ is disjoint from $\Gamma (\mathrm{cl}(MB))$ and $\Gamma(\mathrm{cl}(MC))$ except for $\Gamma (M)$. This follows via saturation, but we want a bit more. Now for $b \in \Gamma ( \mathrm{cl}(MB))$, $c \in \Gamma (\mathrm{cl}(MC))$ and $a \in \Gamma (\mathrm{cl}(MA_1))$ we also demand that
$$E(\sigma ( a ), b ) \leftrightarrow E(a, b)$$ and also that $$E(\sigma (a), c ) \leftrightarrow E( f(a) , c ). $$ This all agrees on the portion which overlaps, since $\mathrm{cl}(MC) \ind _M \mathrm{cl}(MB)$ means that the graph portions of $\mathrm{cl}(MC)$ and $\mathrm{cl}(MB)$ are disjoint except for the graph portion of $M$. Then by saturation, we have such a $\sigma.$  

This $\sigma$ induces an embedding of $Q_{\Gamma ( \mathrm{cl}(MA_1))}$ into $Q_{\sigma(\Gamma (\mathrm{cl}((MA_1))))}$. Now we want to extend $\sigma $ to the module sort. All of the bases which have appeared so far are small relative to the saturation of $\m M$, and so we can pick some $B_{A_1}^*$ with the property that $B_M \cup B_0 \cup C_0 \cup B_{A_1}^*$ is independent over $Q_{\Gamma (\m M)}$. 

Now $(\sigma (\Gamma (\mathrm{cl}(MA_1)))),\span _ {Q_{\sigma (\Gamma (\mathrm{cl}(MA_1)))}} (B_M \cup B_{A_1}^*) $
is a strong closed substructure of $\m M$. 
Define $\sigma^*$
a map from $\mathrm{cl}(MA_1)$
to $\left( \sigma (\Gamma (\mathrm{cl}(MA_1))),\span_{Q_{\sigma (\Gamma (MA_1)))}} (B_M \cup B_{A_1}^*) \right)$ via $\sigma $ on the graph sort. On the  module sort, note that $\sigma $ induces an isomorphism of the associated division rings, and we extend by linearity the above given map of the bases. Call the image of $A_1$ under this map $A_1^*.$ 

The closed structure $\left( \sigma (\Gamma (\mathrm{cl}(MA_1)))),\span _ {Q_{\sigma (\Gamma (\mathrm{cl}(MA_1)))}} (B_M \cup B_{A_1}^* ) \right)$ is the closure of $MA_{1}^*.$ By the claim at the beginning of this proof, $\mathrm{cl}(MA_1B) $ has graph sort $\Gamma ( \mathrm{cl}(MA_1)) \cup \Gamma (\mathrm{cl}(MB))$ and the module sort has basis $B_M \cup B_{A_1} \cup B_0.$ 
Similarly, $\mathrm{cl}(MA_1^*B)$ has graph sort $\sigma (\Gamma (\mathrm{cl}(MA_1))) \cup \Gamma (\mathrm{cl}(MB))$ and basis $B_M \cup B_{A_1}^* \cup B_0$. %%The map $\sigma $ extended by the identity on $\Gamma (\mathrm{cl}(MB))$ is a graph isomorphism between the graphs of these two substructures, while the above defined map of bases and its extension by linearity gives an isomorphism of the module sorts. 
Putting these together, $\sigma $ extended by the identity on $\Gamma (\mathrm{cl}(MB))$ is a graph isomorphism between the graphs of these two substructures, while the above defined map of bases and its extension by linearity gives an isomorphism of the module sorts. Therefore, 
$$\mathrm{cl}(MA_1B) \cong _{MB} \mathrm{cl}(MA_1^*B)$$ and now we know by our characterization of types that $A_1 \equiv _{MB} A_1^*$. The analogous argument gives that $A_1^* \equiv _{MC} A_2$. 

The independence of $A_1^*$ from $BC$ over $M$ follows from the fact that $B_M \cup B_0 \cup C_0 \cup B_{A_1}^*$ is independent and the graphs $\Gamma(\mathrm{cl}(MA_1^*)) $ and $\Gamma ( \mathrm{cl}(MBC))$ have intersection $\Gamma (M)$. 
\end{proof}

\end{document}
